% anexos/anexoF_soluciones.tex
\chapter{Soluciones de los ejercicios}\label{cap:soluciones}

Cada solución se presenta con el código que la resuelve y su resultado. El
mismo código está en el notebook de soluciones de cada capítulo.

% ---------------------------------------------------------------------------
\section*{Capítulo \ref{cap:primer_notebook}}
\addcontentsline{toc}{section}{Capítulo \ref{cap:primer_notebook}}

\soluciondel{2.1}
\fragmento{s02_ejercicio_1}
\salida{s02_ejercicio_1}
La barra es horizontal y apunta en el sentido positivo de $x$: su longitud es
la diferencia de abscisas, $c = 1$ y $s = 0$.

\soluciondel{2.2}
\fragmento{s02_ejercicio_2}
\salida{s02_ejercicio_2}

\soluciondel{2.3}
Con $E = \SI{200}{\giga\pascal} = \SI{2.0e8}{\kilo\newton\per\metre\squared}$ y
$A = \SI{20}{\centi\metre\squared} = \SI{0.0020}{\metre\squared}$:
\fragmento{s02_ejercicio_3}
\salida{s02_ejercicio_3}

\soluciondel{2.4}
\fragmento{s02_ejercicio_4}
\salida{s02_ejercicio_4}
La barra apunta hacia la izquierda y hacia arriba, de modo que su ángulo está
en el segundo cuadrante: el correcto es el de \py{math.atan2}. El cociente
$\Delta y / \Delta x = 3/(-6)$ es igual al de una barra que apunta hacia la
derecha y hacia abajo, y \py{math.atan} no puede distinguir ambos casos.

\soluciondel{2.5}
\fragmento{s02_ejercicio_5}
\salida{s02_ejercicio_5}
Con \py{math.atan(dy / dx)} la división por $\Delta x = 0$ produce el error
\py{ZeroDivisionError} y el programa se detiene.
